% CONSORT Randomized Controlled Trial Template
% Based on CONSORT 2010 Statement (25-item checklist)
\documentclass[12pt]{article}
\usepackage{geometry,graphicx,booktabs,hyperref,natbib}
\geometry{margin=1in}

\title{[TITLE]: A Randomized Controlled Trial}
\author{[AUTHORS]}
\date{}

\begin{document}
\maketitle

\begin{abstract}
\textbf{Background:} [Scientific background and rationale]
\textbf{Methods:} [Trial design, participants, interventions, primary outcome]
\textbf{Results:} [Key findings with effect sizes and confidence intervals]
\textbf{Conclusions:} [Interpretation and implications]
\textbf{Trial Registration:} [Registry and number]
\end{abstract}

\section{Introduction}
% CONSORT Item 2a: Scientific background and explanation of rationale
% CONSORT Item 2b: Specific objectives or hypotheses

\section{Methods}

\subsection{Trial Design}
% CONSORT Item 3a: Description of trial design (parallel, factorial, etc.)
% CONSORT Item 3b: Important changes to methods after trial commencement

\subsection{Participants}
% CONSORT Item 4a: Eligibility criteria for participants
% CONSORT Item 4b: Settings and locations where data were collected

\subsection{Interventions}
% CONSORT Item 5: Interventions for each group with sufficient detail

\subsection{Outcomes}
% CONSORT Item 6a: Pre-specified primary and secondary outcome measures
% CONSORT Item 6b: Any changes to trial outcomes after commencement

\subsection{Sample Size}
% CONSORT Item 7a: How sample size was determined
% CONSORT Item 7b: Interim analyses and stopping guidelines

\subsection{Randomisation}
% CONSORT Item 8a: Method used to generate random allocation sequence
% CONSORT Item 8b: Type of randomisation; details of restriction
% CONSORT Item 9: Mechanism for allocation concealment
% CONSORT Item 10: Who generated sequence, enrolled participants, assigned

\subsection{Blinding}
% CONSORT Item 11a: Who was blinded after assignment
% CONSORT Item 11b: Description of similarity of interventions

\subsection{Statistical Methods}
% CONSORT Item 12a: Statistical methods for comparing groups
% CONSORT Item 12b: Methods for subgroup and adjusted analyses

\section{Results}

\subsection{Participant Flow}
% CONSORT Item 13a: Flow of participants through each stage (diagram)
% CONSORT Item 13b: Losses and exclusions after randomisation

\subsection{Recruitment}
% CONSORT Item 14a: Dates defining periods of recruitment and follow-up
% CONSORT Item 14b: Why the trial ended or was stopped

\subsection{Baseline Data}
% CONSORT Item 15: Baseline demographic and clinical characteristics

\subsection{Numbers Analysed}
% CONSORT Item 16: Number of participants in each analysis group

\subsection{Outcomes and Estimation}
% CONSORT Item 17a: Primary outcome results with effect size and precision
% CONSORT Item 17b: Binary outcomes - absolute and relative effect sizes

\subsection{Ancillary Analyses}
% CONSORT Item 18: Results of any other analyses (subgroup, adjusted)

\subsection{Harms}
% CONSORT Item 19: All important harms or unintended effects

\section{Discussion}
% CONSORT Item 20: Trial limitations (sources of bias, imprecision)
% CONSORT Item 21: Generalisability (external validity, applicability)
% CONSORT Item 22: Interpretation consistent with results

\bibliographystyle{vancouver}
\bibliography{references}
\end{document}
